% Web専用原稿．
% Lang, Algebra (2002), Ch. III §6, pp.142–143 の双対空間・評価写像・双対基底を確認．有限次元と一般の場合を区別する．
\section{代数的双対}
% genron:begin K-000649
% genron:use K-000001
% genron:use K-000002
% genron:use K-000003
% genron:use K-000004
% genron:use K-000005
% genron:use K-000007
% genron:use K-000017
% genron:use K-000560
% genron:use K-000561
% genron:use K-000343
% genron:use K-000047
% genron:use K-000562
\begin{defn}[双対空間]
体 $K$ 上の線形空間 $V$ に対し，$\mathcal L_K(V,K)$ を $V$ の双対空間といい，$V^*$ と書く．その元を線形汎函数という．演算は各点における和及びスカラー倍である．零函数が含まれ，任意の二つの線形汎函数の線形結合も線形なので，これらの演算に関し線形空間をなす．
\end{defn}
% genron:end K-000649
% genron:begin K-000650
% genron:use K-000001
% genron:use K-000002
% genron:use K-000003
% genron:use K-000004
% genron:use K-000005
% genron:use K-000007
% genron:use K-000017
% genron:use K-000649
% genron:use K-000561
% genron:use K-000338
% genron:use K-000036
\begin{defn}[双対写像]
線形写像 $T:V\to W$ に対し，$\varphi\in W^*$ を $\varphi\circ T\in V^*$ に移す写像を $T$ の双対写像といい，$T^\vee:W^*\to V^*$ と書く．合成の線形性及び点ごとの演算から $T^\vee$ は線形である．
\end{defn}
% genron:end K-000650
% genron:begin K-000651
% genron:use K-000001
% genron:use K-000002
% genron:use K-000003
% genron:use K-000004
% genron:use K-000005
% genron:use K-000007
% genron:use K-000017
% genron:use K-000650
% genron:use K-000036
\begin{prop}[双対写像と合成]
線形写像 $S:U\to V,\ T:V\to W$ に対して $(T\circ S)^\vee=S^\vee\circ T^\vee$ であり，$(\operatorname{id}_V)^\vee=\operatorname{id}_{V^*}$ である．
\end{prop}
\begin{proof}
$\varphi\in W^*$ に対して $(T\circ S)^\vee(\varphi)=\varphi\circ T\circ S=S^\vee(T^\vee(\varphi))$ となる．恒等写像についても合成の性質から従う．
\end{proof}
% genron:end K-000651
% genron:begin K-000652
% genron:use K-000001
% genron:use K-000002
% genron:use K-000003
% genron:use K-000004
% genron:use K-000005
% genron:use K-000007
% genron:use K-000017
% genron:use K-000649
% genron:use K-000562
\begin{defn}[零化空間]
線形空間 $V$ の部分集合 $S$ に対して
\[
 S^\perp=\{\varphi\in V^*\mid \forall s\in S,\ \varphi(s)=0\}
\]
を $S$ の零化空間という．零函数を含み，和及びスカラー倍について閉じるから部分空間である．
\end{defn}
% genron:end K-000652
% genron:begin K-000653
% genron:use K-000001
% genron:use K-000002
% genron:use K-000003
% genron:use K-000004
% genron:use K-000005
% genron:use K-000007
% genron:use K-000017
% genron:use K-000650
% genron:use K-000652
% genron:use K-000629
% genron:use K-000634
% genron:use K-000631
\begin{prop}[双対写像の核と像]
線形写像 $T:V\to W$ に対して
\[
 \ker T^\vee=(\operatorname{Im}T)^\perp,\qquad
 \operatorname{Im}T^\vee=(\ker T)^\perp
\]
となる．
\end{prop}
\begin{proof}
第一の等式は $\varphi\circ T=0$ が $\varphi$ が $\operatorname{Im}T$ 上で零であることと同値であることによる．第二の等式について，$\varphi\circ T$ は $\ker T$ 上で零である．逆に $\lambda\in(\ker T)^\perp$ ならば $\psi(T(v))=\lambda(v)$ と定めた函数 $\psi:\operatorname{Im}T\to K$ は適法であり線形である．実際，$T(v)=T(v')$ なら $\lambda(v-v')=0$ であり，演算の保存も $T,\lambda$ の線形性から従う．\GenronRef{K-000634}によって $\psi$ を $W$ 上の線形汎函数 $\varphi$ に延長すると $T^\vee(\varphi)=\lambda$ となる．一般の延長には選択公理を用いる．
\end{proof}
% genron:end K-000653
\section{双対基底と二重双対}
% genron:begin K-000654
% genron:use K-000001
% genron:use K-000002
% genron:use K-000003
% genron:use K-000004
% genron:use K-000005
% genron:use K-000007
% genron:use K-000017
% genron:use K-000649
% genron:use K-000633
% genron:use K-000600
% genron:use K-000604
% genron:use K-000567
\begin{prop}[有限次元の双対基底]
有限次元線形空間 $V$ の基底を $(e_1,\ldots,e_n)$ とする．$e^i(e_j)=\delta_{ij}$（$i=j$ なら $1$，それ以外なら $0$）を満たす線形汎函数 $e^i$ が一意に存在する．$(e^1,\ldots,e^n)$ は $V^*$ の基底であり，これを双対基底という．特に $\dim_K V^*=\dim_K V$ である．
\end{prop}
\begin{proof}
\GenronRef{K-000633}により各 $e^i$ が一意に存在する．任意の $\varphi\in V^*$ に対し $\varphi=\sum_i\varphi(e_i)e^i$ が基底上で成立するので全域で成立する．また $\sum_i a_ie^i=0$ を各 $e_j$ に適用すると $a_j=0$ となる．従って双対基底は線形独立かつ生成族である．
\end{proof}
% genron:end K-000654
% genron:begin K-000655
% genron:use K-000001
% genron:use K-000002
% genron:use K-000003
% genron:use K-000004
% genron:use K-000005
% genron:use K-000007
% genron:use K-000017
% genron:use K-000649
% genron:use K-000654
% genron:use K-000633
% genron:use K-000602
% genron:use K-000638
% genron:use K-000629
\begin{prop}[二重双対への標準写像]
$V^{**}=(V^*)^*$ を二重双対空間という．評価による写像
\[
 \iota_V:V\to V^{**},\qquad \iota_V(v)(\varphi)=\varphi(v)
\]
は線形単射である．$V$ が有限次元なら線形同型である．
\end{prop}
\begin{proof}
各 $v$ に対し評価は $\varphi$ について線形であるから $\iota_V(v)\in V^{**}$ であり，$\varphi$ の線形性により $\iota_V$ も線形である．$v\ne0$ なら $\{v\}$ は線形独立であり，\GenronRef{K-000602}により基底に拡張できる．\GenronRef{K-000633}により $v$ を $1$，他の基底元を零に移す汎函数 $\varphi$ が存在する．従って $\iota_V(v)\ne0$ であり，単射である．一般の場合の基底拡張には選択公理を用いる．有限次元の場合は \GenronRef{K-000654}を二度適用して $\dim_K V^{**}=\dim_K V$ を得るから \GenronRef{K-000638}により全射でもある．
\end{proof}
% genron:end K-000655
% genron:begin K-000656
% genron:use K-000001
% genron:use K-000002
% genron:use K-000003
% genron:use K-000004
% genron:use K-000005
% genron:use K-000007
% genron:use K-000017
% genron:use K-000655
% genron:use K-000650
% genron:use K-000036
\begin{prop}[二重双対と線形写像の整合性]
線形写像 $T:V\to W$ に対して
\[
 (T^\vee)^\vee\circ\iota_V=\iota_W\circ T
\]
が成立する．
\end{prop}
\begin{proof}
任意の $v\in V,\ \varphi\in W^*$ について左辺の $v$ における値を $\varphi$ に適用すると $\iota_V(v)(T^\vee(\varphi))=\varphi(T(v))$ となる．これは右辺の対応する値に等しい．
\end{proof}
% genron:end K-000656
